\bookchapter{Глава 23: Программное мышление и введение в автоматизацию}

\chapterintro{В двух предыдущих главах мы выстроили «стену безопасности» для телефона и увидели, как внутри неё работает всё более умный «мозг AI». Теперь давайте разблокировать один из высших потенциалов телефона: заставить его «думать» как программист и автоматически выполнять заданную вами цепочку действий. Вы скажете: «Я не умею писать код, при чём тут я?» Спокойно — в этой главе вам не придётся написать ни одной строки кода. Современные телефоны снизили порог «программирования» до уровня «перетащи и отпусти».}

\booksection{Что такое автоматизация?}

Простыми словами, автоматизация — это:\textbf{заставить телефон при выполнении определённого условия автоматически делать для вас цепочку действий.}Вам не нужно каждый раз нажимать кнопки и не нужно помнить, что делать — телефон сам определит момент и сам выполнит.

Пример: каждое утро вы просыпаетесь и вручную отключаете режим «Не беспокоить», открываете WeChat посмотреть сообщения, открываете приложение погоды узнать, что надеть, открываете NetEase Cloud Music включить песню. Если упаковать эти шаги в автоматический процесс и настроить «каждый день в 7:00» — вы просто берёте телефон, и всё уже готово. Звучит заманчиво?

По сути автоматизация — это логическая цепочка «если… то…», программный if-else. Освоив этот образ мышления, даже если вы никогда не напишете ни строчки кода, ваша системность в решении задач заметно возрастёт. Ключевые элементы:

\begin{enumerate}
  \item \textbf{Условие срабатывания (триггер)}: когда начать? Может быть «каждый день в 7:00», «при подключении к домашнему Wi-Fi», «при прибытии в определённую точку», «при заряде батареи ниже 20\%».
  \item \textbf{Действие}: что конкретно сделать? Отправить сообщение в WeChat, открыть приложение, настроить громкость, сменить обои, включить музыку, показать уведомление… практически всё, что вы делаете вручную, можно автоматизировать.
  \item \textbf{Промежуточная логика}(опционально, но мощно): если ситуация A — делаем X, иначе — Y. Например, «в будний день воспроизвести новостной подкаст, в выходной — музыку». Это называется условным ветвлением.
\end{enumerate}

\booksection{iOS «Быстрые команды» — инструмент автоматизации для пользователей Apple}

Приложение «Быстрые команды» (Shortcuts) у Apple встроенное, скачивать не нужно. Его суть — собирать разные действия вместе, как конструктор. Сначала о ключевых понятиях:

\begin{itemize}
  \item \textbf{«Быстрая команда»}: законченный автоматический процесс; можно запускать вручную или голосом через Siri. Например, вы говорите «Эй Siri, режим дома», и телефон автоматически выполняет заданную цепочку операций.
  \item \textbf{«Автоматизация»}: команда, запускаемая триггером автоматически, без вашего вмешательства. Например, «при подключении к автомобильному Bluetooth автоматически открыть карты». Настраивается на вкладке «Автоматизация» отдельно от «Быстрых команд».
  \item \textbf{«Действие»}: минимальная единица исполнения. Открыть приложение — действие, отправить сообщение — действие, получить погоду — действие, настроить яркость — действие. Вы можете связать сотни действий в одну цепочку.
\end{itemize}

Вот несколько реально полезных примеров — сделайте сами, и всё станет ясно:

\booksubsection{Пример 1: «Утренний ритуал» одним нажатием}

Каждое утро нажимаете кнопку (или говорите Siri «Доброе утро»), и телефон автоматически: получает сегодняшнюю погоду проговаривает погоду и температуру показывает сегодняшние события календаря открывает ваше музыкальное приложение и начинает воспроизведение. Шаги: новая быстрая команда добавить действие «Получить текущую погоду» добавить «Прочитать текст» добавить «Получить события календаря» добавить «Показать напоминание» добавить «Открыть приложение» и выбрать NetEase Cloud Music. Вот так просто — всё перетаскиванием.

\booksubsection{Пример 2: Автоматический «режим энергосбережения»}

На вкладке «Автоматизация» создайте новую выберите «Заряд батареи» установите «ниже 30\%» добавьте действия «Включить режим низкого заряда», «Снизить яркость экрана до 30\%», «Отключить Bluetooth». С тех пор при заряде ниже 30\% телефон автоматически переходит в энергосберегающий режим — без вашего участия.

\booksubsection{Пример 3: Автоматическая запись расходов по сканированию}

Каждый раз после оплаты через WeChat/Alipay телефон считывает уведомление об оплате и автоматически записывает сумму и название магазина в заметку или таблицу Numbers. Для этого используется триггер «Получено SMS/уведомление» в разделе автоматизации (iOS 17+ поддерживает чтение транзакций Wallet) вместе с действием «Добавить в Numbers». Для тех, кто хочет вести учёт, но ленится вручную — спасение.

Ещё один малоизвестный трюк: быстрые команды можно добавить на домашний экран как иконку приложения. Например, создайте команду «Вызвать такси домой», положите иконку на рабочий стол — одним нажатием откроется DiDi с предустановленным адресом дома, быстрее, чем вручную открывать приложение и вводить адрес.

\booksection{Автоматизация на Android — Tasker и его друзья}

На Android всё чуть сложнее, но и потолок выше. Начнём с простого:

\booksubsection{Начальный уровень: Автокликер}

Многие Android-телефоны имеют встроенный «автокликер» или его можно скачать из магазина приложений. Принцип максимально прост: вы один раз выполняете действия вручную, он записывает координаты нажатий и интервалы, а затем может воспроизвести на высокой скорости. Например, в каком-то приложении каждый день нужно забирать бонус за вход — записываете последовательность нажатий, и потом каждый день она запускается сама. Минус: он привязан к координатам, и если интерфейс изменился (например, выскочила реклама), нажатия уйдут мимо.

\booksubsection{Продвинутый уровень: Tasker — король автоматизации Android}

Tasker — признанно сильнейший инструмент автоматизации на Android, бессменный лидер рейтингов. Честно говоря, кривая обучения крутовата, интерфейс не красавец, но как только освоитесь — покажется, что телефон стал всесильным. Ядро Tasker — та же модель «триггер задача»:

\begin{itemize}
  \item \textbf{Профиль (Profile)}: определяет «когда сработать». Триггером может быть время, место, состояние Wi-Fi, подключение Bluetooth, получение определённого уведомления, уровень заряда, ориентация экрана… Tasker поддерживает десятки триггеров.
  \item \textbf{Задача (Task)}: определяет «что делать после срабатывания». Задача содержит несколько действий, выполняемых последовательно. Типы действий: настройка системы (громкость, яркость, Wi-Fi, Bluetooth, авиарежим), отправка SMS, HTTP-запрос, выполнение Shell-команды, управление воспроизведением, показ диалогового окна и др.
\end{itemize}

Несколько классических сценариев Tasker:

\begin{itemize}
  \item \textbf{Автовключение Wi-Fi по приходу домой}: триггер — «подключение к определённой базовой станции» (около дома вы подключаетесь к фиксированной станции) задача «Включить Wi-Fi». GPS не нужен, заряд не тратится — телефон в районе дома сам включает Wi-Fi.
  \item \textbf{Режим совещания}: триггер — «заголовок события календаря содержит 'совещание'» задача «Беззвучный режим, выключить медиа-громкость, автоответ в WeChat (нужен плагин)». Больше не страшно, что на совещании телефон зазвонит.
  \item \textbf{Ночная защита глаз}: триггер — «время 22:00–06:00» задача «Включить режим защиты глаз, снизить яркость до 20\%, перевести экран в ч/б (режим оттенков серого)». Ч/б экран чудодейственно гасит желание скроллить шортсы — картинка сереет, и листать не хочется.
  \item \textbf{Автоподкаст при подключении наушников}: триггер — «подключение проводных/Bluetooth-наушников» задача «Открыть приложение Xiaoyuzhou и воспроизвести последний выпуск». Надели наушники — и уже слушаете, не нужно разблокировать телефон и искать приложение.
\end{itemize}

Если Tasker кажется слишком сложным, попробуйте его упрощённого собрата:\textbf{MacroDroid}. Интерфейс дружелюбнее, шаблонов много — подходит для тех, кто не хочет заморачиваться. Кроме того, многие китайские Android-системы (MIUI от Xiaomi, HarmonyOS от Huawei, ColorOS от OPPO) уже имеют встроенные функции «сценариев» или «умного помощника» — не такие мощные, как Tasker, но для повседневной автоматизации вполне достаточно.

\booksection{Повседневное применение мышления автоматизации}

Главная ценность освоения автоматизации — не экономия нескольких секунд, а\textbf{изменение способа мышления}. Как только привыкнете к модели «если произошло X, автоматически делаю Y», вы обнаружите, что куча повторяющихся дел исчезнет. Например:

\begin{itemize}
  \item Каждый раз после съёмки — автоматически бэкап в облако, а локальные фото сжать до малого размера для экономии места.
  \item Каждое утро автоматически читать расписание из календаря и отправлять текст семье через WeChat.
  \item При получении SMS о доставке — автоматически добавить в заметки задачу «забрать посылку».
  \item При полной зарядке — голосовое напоминание отключить зарядное устройство (предотвращает перезаряд, портящий батарею).
  \item Время окончания рабочего дня — если вы ещё рядом с офисом (по базовой станции) — напоминание «пора домой».
\end{itemize}

Звучит продвинуто, но на деле — как только вы усвоите два базовых понятия «триггер + действие» и соедините их с Быстрыми командами или Tasker, всё это реально сделать. Ключ —\textbf{сначала осознать, что «это можно автоматизировать»}: многие не потому что не умеют, а потому что вообще не задумываются в таком направлении.

\booksection{Руководство по подводным камням: частые проблемы автоматизации}

Автоматизация — не всегда гладкий путь, вот несколько ям, на которые я уже наступил за вас:

\begin{itemize}
  \item \textbf{Геотриггер может не срабатывать}: «Прибытие/убытие» в iOS использует геозабор; точность зависит от сигнала GPS и системной оптимизации. Иногда триггер срабатывает лишь через 10 минут после прихода домой. Лучше использовать подключение к Wi-Fi или базовой станции — надёжнее и экономичнее.
  \item \textbf{«Автоматизации» в Быстрых командах требуют подтверждения}: в iOS большинство триггеров на вкладке «Автоматизация» требуют ручного подтверждения (механизм безопасности). Но триггеры по подключению Wi-Fi, Bluetooth и остановке будильника могут работать полностью автоматически, без подтверждения.
  \item \textbf{Tasker требует разрешения доступности}: многие продвинутые функции (например, имитация нажатий) требуют разрешения службы доступности, из-за чего в настройках висит постоянное уведомление. Это дизайн безопасности Android — убрать нельзя.
  \item \textbf{Не делайте автоматизацию слишком сложной}: начинайте с малых процессов из 2–3 действий. Не бросайтесь сразу строить 20-шаговую суперсхему — отладка сведёт с ума.
\end{itemize}

\textbf{Совет: правильный подход к началу автоматизации}\par

Потратьте 5 минут и подумайте: какие 3 действия вы чаще всего повторяете на телефоне? Ежеутренняя рутина, привычка при подключении Bluetooth-наушников, настройки перед выходом… С этих частых мелочей начните, по одной простой автоматизации на каждую. Даже если каждая экономит лишь 5 секунд, за месяц набежит солидная сумма. А чувство удовлетворения от готовой автоматизации — по-настоящему затягивает.

\begin{tcbpractice}{Попробуйте: три задания для начинающих}
Выберите самое подходящее и выполните — и вы уже в деле:

\begin{enumerate}
  \item \textbf{Задание 1: Утренняя автоматизация}: iOS: в приложении «Быстрые команды» создайте команду добавьте «Получить текущую погоду» «Прочитать текст» «Открыть приложение» (NetEase Cloud Music). После создания скажите Siri «Доброе утро» для запуска. Android: в Tasker создайте задачу добавьте «Получить погоду» «Прочитать» «Открыть приложение» установите триггер на каждый день в 7:00.
  \item \textbf{Задание 2: Автоэнергосбережение при низком заряде}: iOS: вкладка «Автоматизация» новая «Заряд батареи» ниже 25\% добавьте «Включить режим низкого заряда» + «Яркость 25\%». Android: Tasker новый Profile Состояние Питание заряд 0–25\% задача «Включить энергосбережение» + «Снизить яркость».
  \item \textbf{Задание 3: Автонавигация при подключении к автомобильному Bluetooth}: iOS: Автоматизация новая «Bluetooth» выберите свой автомобильный Bluetooth добавьте «Открыть приложение (Amap)» + «Медиа-громкость 80\%». Android: Tasker Profile Состояние Bluetooth подключён выберите автомобильный Bluetooth задача «Открыть Amap» + «Повысить медиа-громкость».
\end{enumerate}

\end{tcbpractice}

\begin{tcbinfo}{Пример 3: Автоматизация «Утреннего ритуала»}
\begin{itemize}
  \item \textbf{Сценарий}: каждый день в 7:00 автоматически проговаривается погода + отключается беззвучный режим + открываются новости.
  \item \textbf{Шаги}: iOS: Быстрые команды Автоматизация триггер по времени добавить «Получить погоду» + «Настроить громкость» + «Открыть приложение»; Android: MacroDroid триггер (время 7:00) действие (медиа-громкость / открыть URL / TTS-прочтение погоды).
  \item \textbf{Ожидаемый результат}: работает без установки, снижая порог входа.
\end{itemize}

\end{tcbinfo}
